\section{Outline of Computational Procedure}

\begin{figure*}
  \centering

  \includegraphics[width=0.70\textwidth]{flowchart.pdf}
  \caption{flowchart for computational procedure }
  \label{fig:flowchart}
\end{figure*}

(flowchart)
The procedure for calculating the piezoelectric matrix (consider renaming) will follow
that presented in \fig{fig:flowchart}.

\subsection{Geometry Optimization}

The first step is of course to optimize the geometry of the molecule/s or small system.  
Depending the the degree of accuracy desired for the calculation, it may be necessary to use
more stringent thresholds for gradient of the energy and especially displacement (since we are mainly
interested in approximating displacement under the application of an applied field).  
Furthermore, if one is performing a density functional theory (DFT) calculation for the energy, it might be necessary
to use a finer grid for the integration and smaller cutoffs for the integrals. 
We have noticed previously that without a decent grid and low cutoffs that there tends to be small ``bumps'' in  cross-sections of the 
energy which may interfere with calcuations of the hessian and other derivatives with respect to displacement.

Furthermore, one might be interested in constraining the system in some position which might be relevant to the crystal 
structure or some other reason.  
For these purposes it might be necessary to employ different methods for constraining the
geometry and still performing the geometry optimization.  
These might include lagrange multiplier or projection techniques.
For these we prefer the reader to \cite{numRec, schlegel, baker}.  
It is important to remember though that in later steps one might want to 
later project out such unwanted degrees of freedom in addition to rotations and translations.
Any type of motion corresponding to an unwanted deformation can be constructed just like the rotation and translation vectors in the Appendix and 
also projected out of the basis.  
The corresponding transformation column vector matrix ${\bf{V}}$ will then be reduced in columns by the number of additional unwanted vectors.  
This would of course not hold if these other unwanted motions are linearly dependent with the rotations and tanslation vectors.
\subsection{Hessian Calculation}

Frequency calcualtions have become statndard in many program packages and often the Hessian matrix is reported at the 
end of the calculation or can be recovered from scratch files.  
Analytical Hessian calculations
might be prohibitively expensive for large systems and parallel implementaions are not always present for some methods in software packages.
However, numerical calculations of the Hessian are always 
embarassingly parallel due to the vast number of independent gradient calculations.  
We have found for our calculations that molecules greater than 25 atoms or so start to become
prohibitively expensive for analytical B3LYP \cite{b3lyp,Lee1988} calculations with a moderately
sized basis set (6-31G(d))\cite{pople}.
Furthermore, we do not notice a remarkable discrepancy in the 
final calculation of the piezoelectric matrix when using a numerical Hessian.  

\subsection{Dipole Derivative Calculation}

The dipole derivative, $\frac{\partial^2 E}{\partial \uu \partial \ff}$, calculation may be performed in parallel with the Hessian calculation.
Alternatively, if one is performing a vibrational analysis for the Hessian calculation, the 
dipole derivatives are usually calculated as well since they have a well established 
relationship with the intensity of infrared active vibrational modes \cite{kom}.  
If this is the case, it is likely that one can merely retrieve the dipole derivatives with the Hessian 
at the end of a frequency calculation. 

If one wishes to perform the dipole derivative calculation manually we refer the reader
to Ref \cite{kom}, and we also present the following procedure.  
After the geometry optimization, either retrieve or perform a separate calculation to obtain the gradient ($\nabla E(\ff=0)$) of the energy. 
After this, one can calculate the gradient of the molecule for three separate field calculations.  
The most reasonable choices for field directions are the x, y and z direction.  
The field magnitude, which we call here $df$,  applied in each direction should be quite
small (~0.001V/nm).  
After the three calculations are finished, one may approximate $\frac{\partial \nabla E}{\partial f_x}$,$\frac{\partial \nabla E}{\partial f_y}$ and 
$\frac{\partial \nabla E}{\partial f_z}$, (or the corresponding length variables for whatever basis you choose) using the finite differences of the gradient from zero field.
\begin{equation} 
\begin{split}
\frac{\partial \nabla E}{\partial f_x} &\approx \frac{\nabla E(f_x = df, f_y=0, f_z=0)-\nabla E(\ff=0)}{df} \\
\frac{\partial \nabla E}{\partial f_y} &\approx \frac{\nabla E(f_x = 0, f_y=df, f_z=0)-\nabla E(\ff=0)}{df} \\
\frac{\partial \nabla E}{\partial f_z} &\approx \frac{\nabla E(f_x = 0, f_y=0, f_z=df)-\nabla E(\ff=0)}{df}
\end{split}
\end{equation}

Each of these derivatives is a vector quantitiy of dimension $3N$, where $N$ is the number of nuclei, and represent the columns of the matrix ${\bf{H_uf}}$ as presented in the previous sections 
and (flowchart figure).
Alternatively, one might choose to perform a few gradient calculations at a few different small field strengths for each field direction and perform a linear regression on each component of the gradient
to find the slope of each component. 
We did not notice that this drastically improves the result and is usually unnecessary.  (consider adding the graphs for this)

\subsection{Project Out Rotations and Translations}

Now that we have the Hessian ${\bf{H_{uu}}}$ and the Dipole derivative ${\bf{H_{uf}}}$ matrices, we need to project out rotations, translations and any other unwanted modes
(if using constraints)from the matrices.  
This is equivalent to transforming the the basis of the displacement variables to a subspace of modes which lack displacement corresponding to rotation
and translation.  
A typical unitary transformation of the Hessian for example (${\bf{H_{UU}}}={\bf{U}}^T \cdot {\bf{H_{uu}}} \cdot {\bf{U}}$) would transform the Hessian from derivatives
$ H_{u_iu_j} = \frac{\partial^2 E}{\partial u_i \partial u=j}$ to $H_{U_iU_j}=\frac{\partial^2E}{\partial U_i \partial U_j}$, where the variables of the transformed matrix are now with respect
to the length variables of the new basis.  
Here we consider that ${\bf{U}}$ is unitary and made up of orthonormal column vectors.  
In the case where we ignore rotations and translations, we must construct the matrix ${\bf{V}}$ discussed in previous sections by projecting out translation and rotation
vectors from the eigenvectors of the the Hessian (or any other basis that spans the full nuclear space).  

  In practice, if one is using modern quantum chemistry software, it might be possible to obtain cartesian normal modes from job outputs or the scratch directory.  
 If this is the case the vectors may already be sorted so that the vectors corresponding to rotations and translations are reported first or last.  
 If this is not the case, one may always apply the Hessian to 
the normal modes, and record the norm of the results.  
The vectors corresponding to the smallest norms (ie. very close to zero) are typically the vectors which represent translations and rotations.
If a molecule is linear we expect there to be 5 such modes, and ${\bf{V}}$ will have dimesions $3N\times3N-5$.  
If the molecule is nonlinear we expect 6, and ${\bf{V}}$ will have dimesions $3N\times3N-6$.
If the translational modes and rotational modes have been identified in this way, the remaining normal modes are make up the columns of ${\bf{V}}$. 
One should also check that these vectors are orthonormal.  
Using ${\bf{V}}$, we may now transform ${\bf{H_{uu}}}$ and ${\bf{H_{uf}}}$ to ${\bf{H_{vv}}}$ and ${\bf{H_{vf}}}$ according to (flowchart and equation).

If the program package one is using does not offer cartesian normal modes after a frequency calculation, it is possible to construct translational and rotational modes which may then 
be projected out from the eigenvectors of the Hessian or mass weighted Hessian.  
This should then lead to $3N-5$ or $3N-6$ (for linear or nonlinear molecules respectively) nonzero 
vectors which may then be internally orthonormalized via the Gram-Schmidt method or otherwise \cite{numRec}.
It might be the case that after projecting out the rotations and translations that there are not $5$ or $6$ zero vectors.
After the orthogonalization routine, there should be this many zero vectors.  
For more information, we refer the reader to Ref \cite{numRec}.

\subsection{Solve for Displacement Derivative Matrix}

Now that we have the transformed Hessian and dipole derivative matrix (${\bf{H_vv}}$ and ${\bf{H_vf}}$ respectively), we may now solve for the displacement derivative matrix 
$\firstDeriv{\uu_{\mathrm{vib}}}{\ff}$.  
This matrix contains the derivatives of all the nuclear displacement vectors ${\bf{u}}$ with respect to the field vector ${\bf{f}}$ under condition
that the energy remains minimized and only deformations in the space of the vibrational motion of the molecule is allowed.  
The matrix is inherently $3Nx3$ once put back into normal cartesian space.  
We generally solve in the vibrational mode space, but simultaneously 
transform back to the cartesian space.  
The equation is given by \eq{eq:finalDxdf} and is also included in \fig{fig:flowchart}.  
Even though the inverse of the ${\bf{H_vv}}$ is written in \eq{eq:finalDxdf}, 
it is not necessary to invert this matrix; it is generally more efficient to just solve equation (equation) for $\frac{\partial {\bf{v}}}{\partial {\bf{f}}}$ via known algorithms, widely implemented in such programs as Mathematica, Matlab, and Python
packages \cite{math,MATLAB:2010,python}.  
Then of course one must transform $\frac{\partial {\bf{v}}}{\partial {\bf{f}}}$ back to normal space via ${\bf{V}}\cdot \frac{\partial {\bf{v}}}{\partial {\bf{f}}}$.

\subsection{Construct Piezoelectric Matrix}

As discussed in the previous Derivation section, to construct the final piezoelectric matrix ${\bf{P}}$, one first needs to define a ``material line'', 
which in the molecular sense we 
take to mean the line segment connecting two atoms of choice in the system (or some other points of reference which will change upon deformation).  
The unit vector in our position coordinate
system then becomes ${\bf{e}}$.  
From the matrix ${\bf{P}}$ we can describe the relative deformation of these two points.  
${\bf{P}}$, to reiterate, will be an approximation to
$\frac{\partial^2 {\bf{u}}}{\partial r_0 \partial {\bf{f}}}$, which we approximate via the change in the derivative of the displacement 
vector ${\bf{u}}$ with respect to the field vector
${\bf{f}}$ as we move from atom ``a'' to atom ``b'', per unit distance as we move from one atom to another along the ``material line'' in the undeformed body.
The calculation is simple and is given by
\eq{eq:INeedIt2} and \fig{fig:flowchart}.  
The only difficult task is identifying which parts of the $3Nx3$ displacement derivative matrix $\frac{\partial {\bf{u}}}{\partial {\bf{f}}}$ correspond to 
$\frac{\partial {\bf{u}}_a}{\partial {\bf{f}}}$ and $\frac{\partial {\bf{u}}_b}{\partial {\bf{f}}}$.  
Both parts are $3x3$ and hence make ${\bf{P}}$ also $3x3$.  

Finding these parts, however, depends on the order of position variables used by the software to write the Hessian and gradients of the molecule.  
Typically the ordering is 
(atom1 x, atom1 y, atom1 z  . . . atomN x, atomN y, atom N z). 
So if one uses for example atom 3 for ``a'' and atom 35 for ``b'' as the two ends of the ``material line'',  $\frac{\partial {\bf{u}}_a}{\partial {\bf{f}}}$
will be given by rows $3*2+1=7$ through $3*2+3=9$, and $\frac{\partial {\bf{u}}_b}{\partial {\bf{f}}}$ will be given by $3*34+1=103$ through $3*34+3=105$ of $\frac{\partial {\bf{u}}}{\partial {\bf{f}}}$.  
Hence, $\frac{\partial {\bf{u}}_i}{\partial {\bf{f}}}$
corresponds to rows $3*(i-1)+1=3i-2$ through $3*(i-1)+3=3i$. 
This of course assumes that we index the rows starting at $1$ and up to $3N$.  
In python and other languages it is common
to index starting from $0$ to $3N-1$. 
If this is the case, then atom 3 would actually be atom 2 (atom 0 and atom 1 come before) and atom 35 would be atom 34.  
In this instance the 
rows to extract from $\frac{\partial {\bf{u}}}{\partial {\bf{f}}}$ would be $3i$ to $3i+2$.  








